\chapter{Scalability}
\label{ch:scalability}

\section{Scope of the question}

Expo Snips is a developer tool distributed to individual workstations and an
editor. It is not a multi-tenant online service, so conventional
\emph{horizontal} web-scale concerns (load balancing, database sharding, request
fan-out) are \textbf{NOT APPLICABLE}. Scalability here means three concrete
things:

\begin{enumerate}
  \item \textbf{Catalog growth} --- does the product still work as the number of
        recipes and templates grows?
  \item \textbf{Consumer base growth} --- does the delivery mechanism (npm,
        GitHub Pages, VS Code Marketplace) absorb many users?
  \item \textbf{Team growth} --- can more contributors add recipes without
        conflicts?
\end{enumerate}

\section{Catalog growth}

\begin{longtable}{p{3.8cm}p{1.9cm}p{8.5cm}}
\caption{Catalog growth analysis}\\
\toprule
\textbf{Dimension} & \textbf{Assessment} & \textbf{Reasoning} \\
\midrule

Lookup & Good & In-memory array; negligible at hundreds of entries. \\
Build time & Good & Linear in recipe count; \texttt{build:registry} walks the
data directory once. \\
Offline bundle size & Concern & Each client inlines the whole catalog. At 24
recipes the CLI's catalog module is $\approx$240\,KB; it grows linearly, so a
catalog of several hundred recipes would bloat the npm package and the extension
bundle proportionally. \\
Template duplication & Concern & shadcn and per-styling overrides duplicate
template text per recipe rather than sharing components. \\
Delivery size limits & Risk & npm enforces package size guidance; very large
catalogs would need to move to the remote-registry path by default. \\
\bottomrule
\end{longtable}

\textbf{INFERRED:} the current architecture scales comfortably to a few dozen
recipes but the inline-everything decision becomes the limiting factor at
scale. The product already contains the escape hatch (the remote registry and
\texttt{--registry}/\texttt{registryUrl}), but it is off by default.

\textbf{RECOMMENDED:} if the catalog is expected to exceed roughly one hundred
recipes, move the default to a lighter client that fetches an index and loads
individual recipes on demand, and keep the full bundle only as an offline
fallback.

\section{Consumer base growth}

\begin{itemize}
  \item \textbf{npm distribution:} the public registry is designed for very high
        read concurrency; the two published packages (\texttt{@expo-snips/core},
        \texttt{@expo-snips/cli}) impose no server-side burden on the project.
  \item \textbf{Registry hosting:} the registry is a set of static JSON files
        deployed to GitHub Pages, which is a CDN-backed static host. Serving
        \texttt{index.json} and per-recipe documents scales without any
        application server. There is no dynamic backend to saturate.
  \item \textbf{VS Code extension:} distribution and updates are handled entirely
        by the Marketplace; the extension makes no call to project-owned
        infrastructure unless the user configures a registry URL.
  \item \textbf{No telemetry:} because nothing reports back, traffic growth has no
        analytics or write path to contend with.
\end{itemize}

\textbf{Status:} IMPLEMENTED. The product's delivery model is effectively
infinitely scalable for its purpose because it is static and distributed.

\section{Team growth}

\begin{itemize}
  \item Recipes are independent JSON documents under
        \texttt{packages/registry/data/}; two contributors adding different
        recipes touch different files, which is a low-conflict shape for git.
  \item The \texttt{validate} script enforces structure per recipe, and the
        \texttt{build:registry} script fails on a missing template, giving
        contributors fast feedback.
  \item There is no shared registry lock file or generated index committed per
        change beyond the build output, so merge conflicts on the index are
        avoided.
\end{itemize}

\textbf{Status:} IMPLEMENTED. \textbf{Gap:} the absence of automated tests and
the absence of a root \texttt{README} or contribution guide make onboarding
rely on reading source (see Chapter~\ref{ch:technical-debt}).

\section{Known scalability limits}

\begin{longtable}{p{1.5cm}p{9.3cm}p{3.3cm}}
\caption{Scalability limits and risks}\\
\toprule
\textbf{ID} & \textbf{Limit} & \textbf{Trigger} \\
\midrule

SCALE-01 & Entire catalog inlined into every client bundle. & Catalog grows past
a few dozen recipes. \\
SCALE-02 & No pagination or search index for the remote registry; the client
fetches and filters the whole index. & Very large catalogs over a remote
registry. \\
SCALE-03 & Single static registry host with no versioning strategy documented for
breaking recipe changes. & Multiple concurrent catalog versions in the wild. \\
SCALE-04 & Dependency installation is serial and user-visible. & Recipes with many
or heavy dependencies. \\
\bottomrule
\end{longtable}

\section{Summary}

\begin{table}[H]
\centering
\caption{Scalability posture}
\begin{tabular}{p{4.4cm}p{9.8cm}}
\toprule
\textbf{Axis} & \textbf{Verdict} \\
\midrule
Runtime concurrency & NOT APPLICABLE (single-user local tool). \\
Delivery / distribution & IMPLEMENTED and inherently scalable (static CDN + npm +
Marketplace). \\
Catalog growth & INFERRED to be the real constraint; bounded by the inline
bundle. \\
Team growth & IMPLEMENTED (independent recipe files) with onboarding gaps. \\
\bottomrule
\end{tabular}
\end{table}

\textbf{Status:} IMPLEMENTED (delivery) / INFERRED (catalog limit) / RECOMMENDED
(remote-first default at scale).
